\chapter{Integración e ingesta de fuentes de datos}
\label{Appendix:integracion}

La aplicación está diseñada para operar sobre datos provenientes de fuentes
externas heterogéneas, lo que obliga a disponer de una capa de integración
que medie entre cada origen y el modelo relacional descrito en el cuerpo de
la memoria. Este apéndice documenta el conjunto de scripts que conforman
dicha capa, distinguiendo entre la descarga remota, el procesado y la
carga final en PostgreSQL. Se detallan los dos orígenes incorporados a
día de hoy, TMDB y DBLP, y se describe el script de poblado, común a
ambos, que materializa los objetos en las cinco tablas del esquema.

\section{Pipeline general}

La integración de cada fuente sigue un esquema de tres etapas
diferenciadas. La primera etapa, de \emph{extracción}, obtiene los datos
brutos desde el sistema de origen, ya sea mediante una API REST autenticada
o a través de volcados públicos comprimidos. La segunda etapa, de
\emph{procesado}, transforma esos datos en una representación intermedia
uniforme, en formato JSON Lines (un objeto JSON por línea), que codifica
de forma homogénea el modelo común de colecciones, entidades, metadatos,
referencias y recursos. La tercera etapa, de \emph{ingesta}, vuelca ese
JSONL en la base de datos PostgreSQL siguiendo el esquema descrito en el
capítulo de desarrollo.

Esta separación tiene dos consecuencias prácticas. Por un lado, permite
incorporar nuevas fuentes implementando únicamente las dos primeras
etapas, ya que la tercera es agnóstica al origen siempre que el JSONL
generado se ajuste al formato esperado. Por otro, hace que los archivos
intermedios actúen como caché reutilizable, evitando rehacer descargas o
procesados costosos cuando solo se modifica una etapa posterior.

\subsection*{Modelo de representación intermedio}

La pieza que articula esta arquitectura es el formato del JSONL
intermedio. El archivo contiene un objeto JSON por línea (lo que
permite procesarlo en streaming sin cargarlo entero en memoria) y se
estructura en dos partes. La primera línea es un objeto de cabecera con
la lista de colecciones presentes en el dataset (nombre e identificador
de cada una). Las líneas siguientes son entidades, cada una con cinco
campos: \texttt{id} y \texttt{collection\_id} que identifican la
entidad dentro de su colección, \texttt{metadata} con los atributos
filtrables agrupados por clave, \texttt{references} con los vínculos
tipados hacia otras entidades del mismo archivo y \texttt{contents}
con el texto descriptivo y los recursos externos. El
Listado~\ref{lst:jsonlSchema} reproduce un fragmento ilustrativo del
formato.

\begin{lstlisting}[style=jsonStyle, caption=Esquema del JSONL intermedio, label=lst:jsonlSchema]
{"collections": [{"name": "Movies", "id": 1}, {"name": "People", "id": 3}]}
{
  "id": 550,
  "collection_id": 1,
  "metadata": {
    "Genres": ["Drama"],
    "Release date (Year)": ["1990-1999"],
    "Runtime": ["120-149"]
  },
  "references": [
    {"reason": "Director", "reference_id": 7467, "reference_collection_id": 3}
  ],
  "contents": {
    "Name": "Fight Club",
    "Overview": "...",
    "_resources": [
      {"type": "image", "label": "Poster", "url": "https://image.tmdb.org/..."}
    ]
  }
}
\end{lstlisting}

Este formato cumple varios objetivos a la vez. Es directamente
proyectable al esquema relacional: cada clave del bloque
\texttt{metadata} produce una o más filas en la tabla \texttt{metadata},
cada elemento de \texttt{references} se convierte en una fila de
\texttt{reference}, y el bloque \texttt{contents} se almacena tal cual
en la columna \texttt{JSONB} de \texttt{entities}. Es también
extensible, ya que la ausencia de un esquema fijo en los bloques
\texttt{metadata} y \texttt{contents} permite que cada colección
exponga los campos propios de su dominio sin necesidad de coordinarse
con las demás. Y, por último, es independiente de la fuente: tanto el
procesado de TMDB como el de DBLP producen archivos con esta misma
estructura, lo que hace que la etapa de ingesta sea un único script
común a todos los orígenes.

\section{TMDB}

The Movie Database (TMDB) es una base de datos colaborativa de
información cinematográfica y televisiva, mantenida por una comunidad
amplia de contribuyentes y consumida por servicios de terceros a través
de una API REST pública. Su catálogo incluye decenas de miles de
películas, series, personas (actores, directores, equipo técnico),
productoras y cadenas de televisión, todos enlazados entre sí por
relaciones tipadas. Esa combinación de tamaño, riqueza de metadatos y
densidad de relaciones la convertía en el banco de pruebas ideal para
un sistema cuyo objetivo es precisamente explorar grafos de información
mediante metadatos y referencias.

La extracción se apoya en un par de scripts en Python independientes,
uno por etapa. El primero descarga los identificadores diarios
publicados por TMDB y, a continuación, consulta el endpoint de detalle
correspondiente para cada uno, escribiendo el JSON devuelto en un
archivo JSONL por colección. El acceso a la API requiere una clave
personal, gestionada como variable de entorno \texttt{TMDB\_API\_KEY}
definida en \texttt{.env.thirdparty} para mantenerla fuera del control
de versiones. La descarga utiliza un pool de corrutinas asíncronas con
un limitador que respeta el ritmo admitido por la API y un mecanismo
de reanudación que evita volver a pedir identificadores ya presentes
en el archivo de salida, lo que hace la descarga idempotente ante
interrupciones.

\subsection*{El interés de trabajar con un subconjunto}

Operar con el catálogo completo durante el desarrollo no era práctico.
Las decenas de miles de entidades por colección, multiplicadas por las
referencias cruzadas entre ellas, habrían provocado tiempos de poblado
de la base de datos y de respuesta del frontend prohibitivos para una
iteración rápida. Más relevante todavía, una parte significativa del
catálogo corresponde a contenidos poco conocidos o con muy pocos votos,
cuya presencia en la exploración aporta ruido sin mejorar el valor
demostrativo del sistema.

La solución adoptada fue puntuar cada entidad mediante un promedio
bayesiano sobre el número de votos y la nota media, y retener
únicamente las primeras \texttt{TOP\_K} entidades de cada colección.
Con \texttt{TOP\_K = 500}, el dataset resultante (denominado
\texttt{dataset\_top500.jsonl}) cabe holgadamente en memoria, se ingesta
en pocos minutos y mantiene precisamente los títulos sobre los que el
usuario tiene más probabilidades de querer navegar. El mismo script
admite ampliar el subconjunto a varios miles de entidades sin tocar
código, ajustando únicamente el parámetro de corte, lo que permite
escalar la prueba con escaso esfuerzo.

\subsection*{Configuración por colección}

El procesado se parametriza mediante \texttt{scripts/TMDB/format.json},
un archivo declarativo que describe las cinco colecciones consideradas
(películas, series, personas, compañías y cadenas) y, para cada una,
los campos que deben extraerse del JSON crudo y cómo se clasifican
dentro del modelo. La configuración define qué campos son
\emph{metadatos} (filtrables), \emph{contenidos} (texto descriptivo) o
\emph{recursos} (URLs externas, como pósteres o enlaces a IMDB), así
como el tipo de cada metadato: cadena, numérico, fecha, ubicación o
codificado. La discretización en rangos para los campos numéricos y de
fechas, motivada en el capítulo de desarrollo, se controla también
desde este archivo mediante un tamaño de bucket por campo.

Centralizar la configuración en este archivo permite extender el
dataset con nuevos campos sin tocar el código Python, y facilita además
que el mismo motor de procesado se aplique a fuentes distintas con un
esquema diferente pero un objetivo análogo, idea que se materializa en
la integración con DBLP.

\section{DBLP}

La segunda fuente integrada es DBLP, un servicio bibliográfico mantenido
por la Universidad de Tréveris que cataloga la producción científica en
informática: artículos en revistas y conferencias, monografías, autores
y lugares de publicación. A diferencia de TMDB, DBLP no expone una API
de consulta detallada; en su lugar, publica un volcado RDF completo en
formato \texttt{N-Triples} (\texttt{dblp.nt.gz}) que, en su variante
íntegra, ronda decenas de gigabytes una vez descomprimido. La fuente es
interesante para este proyecto por dos motivos: aporta un dominio
ortogonal al cinematográfico, lo que valida la generalidad del modelo,
y permite construir un dataset alineado con la comunidad universitaria
filtrando por afiliación.

El procesado se ha implementado en Go por razones de rendimiento: la
lectura completa del volcado debe poder hacerse en tiempos razonables y
con un consumo de memoria acotado. La estrategia se basa en una
primera pasada lineal sobre todas las tripletas que asigna
identificadores enteros a las entidades reconocidas y, en paralelo,
reparte las tripletas en sesenta y cuatro archivos de partición usando
una función hash sobre el sujeto. Esto garantiza que toda la
información de una misma entidad acabe en la misma partición y permite
procesar el grafo en bloques que sí caben en memoria, sin necesidad de
ordenar el archivo completo. Una segunda pasada recorre las
particiones, construye los objetos JSONL aplicando la misma estructura
declarativa de \texttt{format.json} que TMDB y elimina cada partición
inmediatamente después de procesarla para acotar el uso de disco.

Un detalle relevante es el filtrado opcional por afiliación a la UCM:
cuando está activo, el script conserva únicamente personas cuya
información contiene alguna referencia a \texttt{.ucm.es} y, en una
fase de limpieza posterior, descarta las obras que ya no quedan ligadas
a ninguna persona conservada. Este filtro reduce drásticamente el
tamaño del dataset y produce un grafo autoconsistente alineado con el
caso de uso académico previsto. La fase de limpieza realiza además una
poda de referencias colgantes (apuntando a entidades que dejaron de
formar parte del subconjunto), de modo que el JSONL final puede
ingerirse directamente sin pasos de validación adicionales.

\section{Ingesta en PostgreSQL}

La etapa final del pipeline es común a ambas fuentes y se implementa en
\texttt{scripts/populate\_db\_jsonl.py}. El script consume un archivo
JSONL cuyo primer registro describe las colecciones y los registros
siguientes son objetos individuales, y los inserta en las cinco tablas
del esquema (\texttt{datasets}, \texttt{collections}, \texttt{entities},
\texttt{metadata} y \texttt{reference}).

La carga se realiza en dos pasadas. La primera pasada inserta las
entidades y construye en memoria un diccionario que asocia los
identificadores originales del JSON con los identificadores internos
asignados por PostgreSQL. Esta tabla de equivalencias es la que permite,
en la segunda pasada, traducir las referencias del archivo en referencias
internas entre filas de la tabla \texttt{entities}, evitando así las
referencias colgantes y garantizando la integridad del grafo. Durante la
carga, las inserciones se agrupan en lotes para reducir el número de
viajes a la base de datos, y los índices secundarios se desactivan
temporalmente y se reconstruyen al finalizar, lo que acelera de forma
notable el poblado inicial.

La conexión a la base de datos se configura íntegramente mediante las
variables de entorno \texttt{DB\_HOST}, \texttt{DB\_NAME},
\texttt{DB\_USER}, \texttt{DB\_PASS} y \texttt{DB\_PORT}, definidas en
\texttt{.env.db}. El script invoca \texttt{psycopg2} sin credenciales
embebidas, lo que mantiene el archivo libre de información sensible y
facilita reutilizarlo entre entornos de desarrollo y producción. El
nombre del dataset y la ruta al JSONL pueden pasarse como argumentos
de la línea de comandos, lo que permite cargar varios datasets
independientes sobre la misma base de datos.

\section{Ejecución completa desde cero}

Para reproducir el pipeline completo desde un repositorio recién
clonado los pasos son los descritos en el Listado~\ref{lst:pipelineRun}.
Se asume que la base de datos ya está arrancada (véase el
Apéndice~\ref{Appendix:despliegue}) y que las variables de entorno
están definidas en \texttt{.env.db} y \texttt{.env.thirdparty}.

\begin{lstlisting}[style=bashStyle, caption=Ejecución del pipeline completo, label=lst:pipelineRun]
# 1. Instalar dependencias Python
cd scripts
pip install -r requirements.txt

# 2. TMDB: descargar y procesar
python TMDB/pull.py
python TMDB/process.py

# 3. (Alternativa) DBLP: procesar el volcado
cd DBLP
go run main.go dblp.nt.gz format.json
cd ..

# 4. Ingesta en PostgreSQL
python populate_db_jsonl.py "./TMDB/dataset_top500.jsonl" "TMDb Top 500"
\end{lstlisting}
